\noindent\textbf{Dificultad:} --

\subsection{Descripción del nivel}
Reutiliza el usuario del nivel 1 y pide un serial con formato \texttt{XXXX-XXXX-XXXX-YYYY}. Hay que entender cómo se generan los tres primeros bloques y qué representa el bloque final \texttt{YYYY}.

\subsection{Análisis estático}
% Cómo se ubicó la rutina de validación: cadenas, referencias cruzadas, funciones relevantes.
\textit{[A completar.]}

\subsection{Análisis dinámico}
% Breakpoints usados, registros/memoria observados en la comparación.
\textit{[A completar.]}

\subsection{Algoritmo de validación}
% Pseudocódigo del algoritmo reconstruido.
\textit{[A completar.]}

\subsection{Solución (keygen)}
\textit{[A completar.]}

\begin{lstlisting}[language=Python]
# Keygen - Nivel 2
# TODO: pegar el código del keygen
\end{lstlisting}

\subsection{Resultado}
Credencial generada y aceptada por el binario (ver Figura~\ref{fig:levels}):

\begin{lstlisting}[language={}]
Username: mati (from Level 1)
Serial: E834-69B6-1B27-9AA5
[+] Level 2 PASSED
\end{lstlisting}
